% docs/style/uni-verzeichnisstyle.tex -- local copy of tex-skripten/formatierungen/uni-verzeichnisstyle.tex
% Copied 2026-10-12 so that the repository compiles stand-alone (GitHub).
% Local version: 1.0.0 -- unchanged copy

% formatierungen/uni-verzeichnisstyle.tex
% Version 1.1.4
% Änderung gegenüber 1.1.3: Spalten nutzen >{\raggedright\arraybackslash}p{..}
% statt reinem p{..}, damit der Zellinhalt oben (statt vertikal zentriert)
% beginnt – sonst verschieben sich einzeilige Einträge optisch nach unten,
% sobald eine benachbarte Spalte (z.B. "Name") zweizeilig wird.

\newglossarystyle{formelverzeichnis}{
	\setglossarystyle{long3col} % Basis-Style
	\renewenvironment{theglossary}
	{\begin{longtable}{
				@{}
				>{\raggedright\arraybackslash}p{2cm}
				>{\raggedright\arraybackslash}p{2.5cm}
				>{\raggedright\arraybackslash}p{3cm}
				>{\raggedright\arraybackslash}p{\dimexpr\linewidth-7.5cm-4\tabcolsep\relax}
		}}
		{\end{longtable}}
	\renewcommand*{\glossaryheader}{
		\textbf{Symbol} & \textbf{Einheit} & \textbf{Name} & \textbf{Beschreibung} \\ \endhead}
	\renewcommand*{\glossentry}[2]{
		\glsentrysymbol{##1} & \glsentryunit{##1} & \glstarget{##1}{\glossentryname{##1}} & \glossentrydesc{##1} \\
	}
}